\section{Introduzione e obiettivi}

L'algoritmo di Shor consente di fattorizzare numeri interi in tempo polinomiale rispetto
alla lunghezza binaria dell'input, con implicazioni per la sicurezza dei sistemi
crittografici fondati sulla difficoltà della fattorizzazione
\cite{shor1994,nist2024pqc,gidney2021}. La
realizzazione di questo vantaggio teorico richiede tuttavia un controllo degli errori
compatibile con la profondità dei circuiti. Errori nelle operazioni, decoerenza e
lettura imperfetta possono alterare la distribuzione delle misure e compromettere
l'informazione necessaria alla ricerca dell'ordine moltiplicativo \cite{preskill2018}.

La tesi sviluppa uno studio simulativo volto a stabilire quali interventi migliorino
effettivamente l'affidabilità della computazione e rispetto a quali riferimenti. Il
percorso comprende due fasi principali: la valutazione del post-processing classico
degli esiti di Shor, con uno studio di ablazione del contributo del machine learning, e
la caratterizzazione di codici di correzione quantistica e decoder delle sindromi.
Completano il lavoro un'analisi separata della sensibilità di Shor agli errori di porta
e una campagna sull'ottimizzazione circuitale mediante trasformata di Fourier
quantistica approssimata, o AQFT.

I tre ambiti di intervento --- selezione degli esiti, protezione dell'informazione e
riduzione del costo circuitale --- sono collegati da un criterio comune: confrontare
ogni soluzione con un riferimento esplicito, isolare il contributo dei singoli
componenti e misurare il beneficio mediante indicatori coerenti con il compito. Le
campagne rimangono distinte: non costituiscono un'esecuzione integrata di Shor protetta
da correzione d'errore.
